% ============================================================
%  三、求解与结果
% ============================================================
\section{求解与结果}

\subsection{时间序列分析}

图~\cref{fig:timeseries} 展示合成周度指标；亦可直接用 \texttt{pgfplots} 读取 CSV：

\begin{figure}[H]
  \centering
  \begin{tikzpicture}
    \begin{axis}[
      width=0.88\textwidth,
      height=5.2cm,
      xlabel={周次序号},
      ylabel={指标值},
      grid=major,
      legend pos=north west,
      font=\small,
    ]
      \addplot[blue, thick] table[
        x expr=\coordindex,
        y=value,
        col sep=comma,
      ] {../dummy_data/time_series.csv};
      \addlegendentry{观测值}
      \addplot[gray, dashed] table[
        x expr=\coordindex,
        y=trend,
        col sep=comma,
      ] {../dummy_data/time_series.csv};
      \addlegendentry{趋势项}
    \end{axis}
  \end{tikzpicture}
  \caption{由 CSV 直接绘制的时序曲线（与 \cref{fig:timeseries} 数据源一致）}
  \label{fig:ts-pgf}
\end{figure}

\figplot{合成时间序列曲线}{5.8cm}{fig:timeseries}{%
  Matplotlib 自 CSV 导出}{fig_timeseries.pdf}{%
  数据文件：\texttt{time\_series.csv}。}

\subsection{分类结果}

\figplot{二维分类散点图}{6.2cm}{fig:classification}{%
  类别 0/1 散点对比}{fig_classification.pdf}{%
  数据文件：\texttt{classification.csv}。}

\subsection{回归样本表}

表~\cref{tab:reg-sample} 自 CSV 导入前 6 行；完整数据见 \texttt{regression.csv}。

\begin{table}[H]
  \centering
  \caption{回归数据片段（csvsimple 导入）}\label{tab:reg-sample}
  \begin{tabular}{rrrr}
    \topline
    feature & strength & noise & target \\
    \midline
    \csvreader[
      head to column names,
      filter=\ifnum\csvinputline<7,
      late after line=\\,
    ]{../dummy_data/regression.csv}{%
      \feature=\feature, \strength=\strength, \noise=\noise, \target=\target
    }{%
      \feature & \strength & \noise & \target
    }\\
    \bottomline
  \end{tabular}
\end{table}

对比 \cref{fig:ts-pgf}、\cref{fig:timeseries} 与表~\cref{tab:reg-sample} 可见图文混排与智能引用效果。
